\section{Related Work}
\label{sec:related_work}

Connected displays have been developed for activity timing and information sharing in public spaces. Kumar et al.~\cite{kumar2023noticeboard} developed a web controlled noticeboard using an ESP8266-based NodeMCU and 8$\times$8 LED matrix displays. Their interface supports showing messages and controlling scrolling speed and direction. Our system uses two chained 32$\times$16 P10 display large enough for distant viewing, and supports more features including everyday timekeeping, stopwatch, and scheduled count-up and countdown functions. It also integrates device-hosted network configuration and signed firmware updates, which are not described in their implementation.

Kader et al.~\cite{kader2019timing} developed an Arduino Uno-based system for managing viva and presentation sessions. It combines a countdown and buzzer alerts with a separate display that shows group numbers and information to waiting candidates. The examiner operates the system through an IR remote, while an RF link connects the display units, with reported operating ranges of 7~m and 11~m, respectively. Our system uses chained HUB12 panels to provide a large display for timekeeping, notices, and scheduled activities, with management through a web dashboard which is accessible from anywhere with internet access.

Sethi et al.~\cite{sethi2015management} developed a cloud-connected signage management system using Linux-based displays, with phone-assisted provisioning and remote administration through SSH and VNC. Our design targets low-cost MCU  with as much control as we can feasibly get. We use an ESP32-S3 running Zephyr RTOS. This avoids requiring a Linux-capable display controller and browser-based rendering on each device.

Memarovic et al.~\cite{memarovic2016networked} reported the development and deployment of a university display network but their work emphasizes stakeholder involvement and deployment. Their study provides context for introducing displays into shared institutional spaces. Our article focuses on the embedded hardware and software that support timekeeping, announcements, scheduled activities, and maintenance, drawing on ten deployed units operated by teachers and staff.

To the best of our knowledge, prior works do not describe a microcontroller-based display that combines (i) everyday timekeeping, scrolling notices and alerts, stopwatch operation, and scheduled count-up and countdown functions; (ii) web-based content and schedule management with local timing execution and persistent configuration; (iii) device-hosted Wi-Fi provisioning with network scanning and management of up to five saved network profiles; and (iv) over-the-air updates using signed firmware. Our implementation integrates these capabilities on an ESP32-S3 running Zephyr RTOS. We additionally contributed the HUB12 display driver to upstream Zephyr, making it reusable beyond this system.

We also designed and built a custom carrier PCB connecting the controller, RTC, display interface, and power stage, along with a separate board for local controls and status indicators.